\documentclass[../main.tex]{subfiles}
\begin{document}
%==========================================TIÊU ĐỀ TẬP========================================
\begin{table}[h]
    \begin{adjustwidth}{-1cm}{}
        \begin{tabular}{>{\centering\arraybackslash}p{6cm}|>{\raggedright\arraybackslash}p{12.5cm}}
            \multirow{5}{6cm}
            {\\[-40pt]
                \flaregothic\fontsize{20pt}{20pt}\selectfont \centering
                \color{eptype}HISTOIRE\color{black}\\
                \freshbot\fontsize{72pt}{72pt}\selectfont
                \color{epnum}19\\[6pt]
                \cafeteria\fontsize{20pt}{20pt}\selectfont\color{epnum}(D-19)\color{black}
                \\[18pt]
                \color{black}
            }
             &
            {\fotskip\fontsize{18pt}{18pt}\selectfont \ruby{防}{ぼう}\ruby{音}{おん}\ruby{室}{しつ}から\ruby{聞}{き}こえる\ruby{歌}{うた}\ruby{声}{ごえ}}
            \\[6pt]
             & {\cafeteria\fontsize{18pt}{18pt}\selectfont The Voice Behind Soundproof Walls} \\[4pt]
             & {\vnmsans\fontsize{18pt}{18pt}\selectfont Giọng hát sau màn hình chuyển động} \\[8pt]
             & {\caviar{\fontsize{14pt}{14pt}\selectfont Nhân vật: \emp{delu} \emp{kuon2} \emp{suzuno2} \emp{kson}}} \\[8pt]
        \end{tabular}
    \end{adjustwidth}
\end{table}
%=========================================TRUYỆN CHÍNH========================================
\color{black}

\txtbx[2]{
    {\color{vol} \uline{Tóm tắt:}} Trước buổi trình diễn 3D sắp tới, Delutaya bước vào một ngày tổng duyệt cùng sân khấu chuyển động đã được dựng sẵn. Cô muốn phần hát và vũ đạo khớp hoàn hảo, nhưng cứ đến đoạn điệp khúc là nhịp thở và bước chân lại lệch nhau. Kson giúp cô nhìn sân khấu từ góc độ người biểu diễn, Kuon và Game tìm ra điểm vào nhạc dễ theo dõi hơn, còn Suzuno mang đồ ăn nhẹ đến đúng lúc. Thay vì dừng lại mỗi khi mắc lỗi, Delutaya học cách tiếp tục tiết mục và hoàn thành một lượt diễn trọn vẹn.
}

\pov{Delutaya}

Một ngày đầu tháng Năm năm 2022, ánh đèn phòng thu tầng hai đã thắp sáng từ lúc trời còn hơi sẫm, trước khi cả căn nhà kịp thức dậy cho bữa sáng. Tôi đã bật dậy lúc năm rưỡi, kéo rèm cửa ra để nhìn những mảng hồng nhạt đầu tiên hiện ra trên chân trời, dù buổi tổng duyệt chính thức chỉ bắt đầu vào chín giờ. Tôi không hề lo lắng đâu, hay ít nhất, đó là câu chuyện tôi tự nhắc nhở chính mình mỗi khi đứng trước gương rửa mặt.

Đó cũng chính là lời lặp lại khi tôi cúi xuống kiểm tra danh sách thiết bị trên bàn lần thứ tư liên tiếp. Tai nghe cách âm màu đen của tôi: đã để sẵn trên giá đỡ, dây được cuốn gọn gàng. Micro condenser yêu thích của tôi, cái mà tôi đã dùng trong ba buổi thu trước đã được gắn chặt vào giá treo, nút kiểm tra pin sáng xanh ổn.

Bản phối mới nhất của bài hát đã lưu cả vào ổ cứng di động và máy tính chính. Đôi giày tập vũ đạo có đế bám sàn tốt nhất của tôi: đã đặt ngay dưới chân bàn, mũi giày chỉ chạm nhẹ vào tấm thảm trải sàn. Cái khăn bông nhỏ dùng để lau mồ hôi, được giặt sạch phơi khô từ hôm qua: gấp thành hình vuông ngay ngắn, đặt cạnh chai nước lọc đã đổ đầy từ sáng sớm.

Tôi ngồi yên một giây, rồi lại nghiêng người nhìn qua danh sách thêm lần nữa, như thể có một mục nào đó sẽ biến mất nếu tôi không chú ý đủ lâu.

``Delutaya-san, cô vừa xác nhận xong tất cả các mục đó chưa đầy sáu mươi giây trước đây thôi,'' giọng đều đặn của Game vang lên qua loa kiểm âm gắn trên trần nhà, cái mà anh ấy luôn bật để có thể nói chuyện với tôi bất cứ lúc nào.

``Tôi biết mà,'' tôi đáp lại, tay vẫn đang lướt nhẹ qua mép giấy in danh sách, ``chỉ là muốn chắc chắn rằng không có gì bị bỏ sót.''

``Tôi đã ghi nhận được sự chắc chắn đó tổng cộng bốn lần kể từ lúc cô bước vào phòng,'' ông ta tiếp tục, xoáy sâu vào sự tỉ mẩn của tôi.

``Vậy thì để lần thứ năm này sẽ làm ông cảm thấy yên tâm hơn hẳn,'' tôi cười nhẹ, cuối cùng cũng rời mắt khỏi cái danh sách.

``Thực ra thì tôi không phải là người đang cần được trấn an ở đây đâu,'' Game đáp lại, và tôi có thể tưởng tượng ra màn hình phụ của ông ấy đang hiển thị dữ liệu nhịp tim của tôi từ chiếc đồng hồ thông minh tôi đeo.

Tôi không trả lời gì nữa, chỉ quay người nhìn về phía màn hình lớn treo ở cuối phòng. Trên đó, mô hình 3D của tôi đang đứng yên lặng giữa sân khấu ảo được dựng hoàn chỉnh, những tia sáng xanh nhạt như sương mai phủ lên một viền sáng mờ quanh đôi vai của mô hình, làm nó trông thật đến lạ.

Hai dải đèn LED nhỏ màu trắng chạy dọc hai bên sàn sân khấu ảo, như những con đường mòn dẫn dắt mắt người xem thẳng đến vị trí trung tâm, nơi mà tôi sẽ đứng trong buổi diễn thực. Toàn bộ sân khấu chuyển động này đã được dựng xong từ tối hôm qua, mọi góc cạnh đã được điều chỉnh, mọi hiệu ứng ánh sáng đã được kiểm tra. Hôm nay, chúng tôi chỉ còn một việc: chạy thử toàn bộ phần trình diễn, từ đầu đến cuối, như thể nó đang diễn ra thật sự.

Không còn gì mới mẻ để phải ra quyết định cuối cùng. Không còn cái tên bài hát nào cần phải đổi ý, không còn bộ trang phục nào phải chọn đi chọn lại, không còn hiệu ứng nào phải chỉnh sửa từ đầu. Tất cả đã được thống nhất, tất cả đã sẵn sàng.

Tôi chỉ cần làm một việc duy nhất: hát trọn vẹn hết bài hát, bước đúng từng vị trí đã được đánh dấu trên sàn, và không tự làm khó bản thân với những yêu cầu quá khắt khe. Tôi kéo chiếc ghế gỗ ra khỏi bàn, đứng dậy thẳng người ở giữa căn phòng rộng rãi. Nhạc nền vẫn chưa được bật lên, toàn bộ phòng thu chỉ còn tiếng kêu nhỏ của máy điều hòa và tiếng thở đều của tôi.

Tôi hít một hơi thật sâu, đầy cả lồng ngực, rồi đếm nhịp một cách chậm rãi trong đầu. Một, hai, ba, bốn. Đến nhịp thứ năm, vai tôi tự nhiên nâng lên một chút, như thể cơ thể đã nhớ được chuyển động đó từ hàng chục lần tập luyện trước. Tôi từ từ thả lỏng chúng xuống, cảm nhận sự căng thẳng nhỏ trong cơ vai tan biến đi.

``Có vẻ như cô đang khởi động cơ thể mà không cần phát nhạc lên,'' Game nhận xét, khi tôi lại lặp lại động tác nâng vai một lần nữa.

``Tôi chỉ đang tập cách lấy hơi đúng nhịp thôi,'' tay tôi đặt nhẹ lên ngực để cảm nhận không khí đi vào và ra khỏi phổi.

``Tôi có thể bật máy đếm nhịp điện tử lên để tiện theo dõi hơn,'' Game đề xuất, tôi biết ông ta đã chuẩn bị sẵn cái đó trong danh sách công cụ của mình.

``Chưa cần đâu,'' tôi lắc đầu, mắt vẫn nhắm lại để tập trung vào hơi thở của mình.

Tôi lại hít vào một lần nữa. Lần này, hơi thở chậm hơn, sâu hơn, lan tỏa khắp toàn bộ cơ thể tôi.

\ct

Anh Kuon bước vào phòng đúng lúc kim đồng hồ chỉ chín giờ kém năm phút, chiếc máy tính bảng kẹp dưới cánh tay còn tập giấy ghi chú đầy những dấu bút màu đỏ xanh cũn cỡn trong tay. Anh dừng chân ở ngưỡng cửa, mắt quét qua căn phòng đang bật đủ đèn rồi dừng lại trên người tôi, giọng trầm ấm vang lên: ``Em đến sớm thế nhỉ?''

``Em muốn ngắm lại sân khấu và kiểm tra vài thứ nhỏ trước khi mọi người bắt đầu chính thức thôi,'' tôi nhấc cằm đáp, tay vẫn đang vô thức xoay vòng vòng trên mép bàn gỗ.

Anh đặt tập giấy dày cộm lên mặt bàn, tiếng giấy khô lướt nhẹ làm tôi quay người nhìn. ``Anh đã in chi tiết lịch trình chạy thử cho hôm nay rồi. Đầu tiên là mười lăm phút kiểm tra toàn bộ hệ thống âm thanh, sau đó mười phút thử nghiệm độ trễ chuyển động của mô hình trên màn hình, xong hết mình nghỉ giải lao mười phút rồi mới chạy thử lượt đầu tiên toàn bộ bài hát.''

Tôi nghiêng người nhìn danh sách được in rõ ràng, từng mục được đánh dấu thứ tự tỉ mỉ, không khỏi hỏi: ``Anh chia nhỏ từng bước ra thế này sao? Em tưởng mình sẽ chạy luôn cả bài từ đầu?''

``Nếu mình vứt tất cả vào cùng một lúc, khi có lỗi xảy ra thì làm sao mình biết vấn đề nằm ở âm thanh hay ở chuyển động được? Phải tách riêng ra thì mới kiểm soát được từng chi tiết,'' Kuon kéo ghế ngồi xuống đối diện tôi, vừa nói vừa mở máy tính bảng của mình ra.

Đúng lúc đó, màn hình phụ phía trên trần nhà bật sáng, giao diện bảng tính trắng hiện ra với những ô trống được đánh dấu sẵn theo từng đoạn bài hát. Game cũng tham gia vào cuộc hội thoại, giọng đều đặn phát ra từ loa: ``Tôi cũng đã tạo bảng đánh dấu điện tử theo từng phần, tương ứng với lịch trình mà Kuon-user vừa nêu. Mọi kết quả kiểm tra sẽ được cập nhật thời gian thực lên đây.''

``Cảm ơn hai người rất nhiều,'' tôi nói thật lòng, tim cảm thấy ấm áp vì hai người đã chuẩn bị chu đáo thay cho tôi.

Anh Kuon dừng tay đang gõ máy tính, hơi nghiêng đầu nhìn tôi một lúc, rồi nhíu mày nói: ``Em có vẻ hơi căng thẳng nhỉ? Lông mày cứ nhíu lại từ nãy giờ.''

``Em đâu có căng thẳng đâu!'' tôi vội vàng phản bác, vừa nói vừa thả lỏng cơ mặt một cách vô thức.

``Thế sao em lại cầm ngược chai nước rồi?'' anh chỉ vào tay tôi, nụ cười nhỏ khẽ nở trên môi.

Tôi vội vàng cúi nhìn xuống bàn tay của mình. Thật vậy, cái nắp chai nhựa đang úp vào lòng bàn tay, phần đáy chai lại hướng ra ngoài, tôi đã cầm ngược từ lúc nào không hay. ``À... đó là vì em vừa định mở nắp nên mới cầm thế thôi,'' tôi vội vàng xoay chai lại đúng chiều, giọng hơi bối rối.

Kuon không tranh luận thêm, chỉ lẳng lặng rút ra một chiếc khăn bông sạch mới từ trong túi, đặt nhẹ lên bàn cạnh chai nước của tôi. ``Mình cứ làm từ từ từng phần một thôi. Không cần phải ép bản thân hoàn hảo ngay từ lượt chạy đầu tiên đâu. Có lỗi thì mình sửa, có gì chưa ổn thì mình điều chỉnh sau.''

Tôi gật đầu đáp, dù trong sâu thẳm lòng mình vẫn muốn chứng minh rằng tôi có thể làm được mọi thứ hoàn hảo từ đầu. Tôi đã tập luyện bài hát này suốt ba tuần liền, mỗi ngày đều dành ra vài tiếng để thuộc lời, thuộc nhịp, thuộc từng bước chân di chuyển trên sân khấu. Trên giấy, trên sơ đồ, mọi thứ đều vừa vặn đến từng milimet, không có một chi tiết nào bị sót. Vậy mà khi đứng trước sân khấu ảo thật sự, trước mô hình 3D của chính mình trên màn hình lớn, lòng tôi lại dấy lên nỗi lo vô hình, muốn kiểm tra lại mọi thứ thêm một lần nữa, chỉ để chắc chắn rằng không có gì có thể đi sai hướng.

``Bắt đầu với phần kiểm tra âm thanh trước nhé,'' anh Kuon phá vỡ khoảng lặng nhỏ, gõ vào mục đầu tiên trên lịch trình.

Game bật bản nhạc nền lên, nhưng chỉ để ở mức âm lượng rất thấp, đủ để chúng tôi nghe rõ các chi tiết nhỏ. Tôi chăm chú lắng nghe đoạn mở đầu của bài hát, rồi giơ tay ra hiệu dừng. Không có tiếng rè nhỏ từ loa, không có âm thanh lạ nào lọt vào, tất cả các nhạc cụ đều được cân bằng tần số đúng y hệt bản phối cuối cùng chúng tôi đã hoàn thành cách đây một tuần.

Tôi nhắm mắt lại, ngân dài một nốt la để kiểm tra độ vang và độ nhạy của micro. Âm thanh truyền vào tai nghe rõ ràng, không bị bể, không bị méo, mọi thứ đều hoạt động hoàn hảo. ``Âm thanh ổn rồi ạ,'' tôi mở mắt nói, thở nhẹ ra một hơi.

Trên màn hình phụ, Game đã tô màu xanh ô đầu tiên của bảng đánh dấu, ghi rõ dòng chữ "Hoàn thành" bên cạnh. Mục đầu tiên trong lịch trình đã được xong xuôi.

\ct

Tiếp theo là giai đoạn thử nghiệm chuyển động của mô hình 3D, bước quan trọng để đảm bảo mọi cử chỉ của tôi trên sàn thực sẽ được phản ánh chính xác lên sân khấu ảo. Tôi bước nhẹ nhàng đến vị trí khởi đầu, được đánh dấu bằng một vòng tròn sơn mờ màu trắng trên sàn phòng thu, đó là điểm căn cứ chính mà mọi chuyển động sẽ bắt đầu. Ngay khi tôi đứng vững, mô hình 3D trên màn hình lớn cũng xuất hiện tại đúng vị trí trung tâm của sân khấu ảo, chỉ chậm hơn một chút nhỏ đến mức gần như không thể nhận ra.

Tôi nghiêng thân mình sang trái một cách từ từ, giữ thẳng lưng để cảm nhận độ nhạy của cảm biến. Mô hình trên màn hình lập tức nghiêng theo, đường nét chuyển động mượt mà không có bất kỳ độ trễ nào. Tôi giơ cao tay phải ra trước mặt, làm động tác chào nhẹ nhàng như khi gặp một người quen. Mô hình cũng đồng thời giơ tay lên, cổ tay xoay đúng góc độ mà tôi vừa làm. Tiếp theo, tôi thực hiện động tác xoay nửa vòng 180 độ, bước chân di chuyển chính xác đến dấu chân thứ hai được vẽ sẵn trên sàn. Khi tôi dừng lại, mô hình cũng đứng yên ngay tại vị trí tương ứng trên sân khấu ảo, không lệch đi một milimet nào.

Anh Kuon cúi nhìn xuống bảng ghi chú mở trên bàn, bút chì di chuyển nhẹ nhàng trên giấy khi kiểm tra từng tiêu chí. ``Độ trễ giữa chuyển động thực và mô hình ổn định hoàn toàn,'' anh lên tiếng để mọi người cùng nghe. ``Tốc độ cập nhật đủ nhanh, không có hiện tượng chuyển động bị giật hay rời rạc.''

Tôi vẫn đứng yên tại chỗ, mắt dán chặt vào màn hình trong lúc nghĩ ngợi. ``Anh có thể tua lại đoạn vừa rồi cho em xem lại được không? Đoạn mà em xoay người đó,'' tôi hỏi, giọng vẫn còn mang chút lo lắng.

Game lập tức thực hiện lệnh, tua lại đoạn ghi hình vừa quay được vài giây trước. Tôi dõi theo từng khung hình khi mô hình thực hiện động tác xoay trên sân khấu ảo. Đường đi của cơ thể trông rất tự nhiên, không có vẻ cứng nhắc hay gượng ép. Tôi kiểm tra kỹ từng chi tiết nhỏ: vạt áo dài của mô hình không bị lỗi xuyên qua thân người, mái tóc dài cũng không bị mắc vào phần cổ áo như đã từng xảy ra trong một vài lần thử trước. Mọi thứ đều hoàn hảo, nhưng tôi vẫn cảm thấy muốn xem lại thêm một lần nữa để thực sự yên tâm.

``Nếu mình cứ tiếp tục tua lại từng động tác nhỏ như thế này, có lẽ cả buổi sẽ trôi qua trước khi chúng ta kịp chạy thử toàn bộ bài hát đâu,'' anh Kuon nhắc nhở một cách nhẹ nhàng, không có ý trách móc gì.

``Em chỉ muốn chắc chắn rằng mọi thứ đều ổn thôi,'' tôi bước trở lại bàn để ngồi xuống. ``Mấy lần trước từng có lỗi đồ họa, nên em hơi\dots''

``Anh hiểu ý của em, nhưng mà mọi thứ trước mắt đều không có vấn đề gì,'' anh ngắt lời, giọng nói chắc chắn. ``Nếu có điểm nào chưa ổn, chúng ta sẽ ghi lại ngay và xử lý nó sau này. Đừng để những chi tiết nhỏ làm chậm lại cả tiến trình.''

Tôi ngước nhìn lên khuôn mặt anh ấy. Anh Kuon hoàn toàn không có vẻ sốt ruột hay khó chịu với sự tỉ mẩn thái quá của tôi. Anh chỉ đơn giản chỉ vào bảng ghi chú trước mặt, nơi tất cả các kết quả kiểm tra từ đầu buổi đã được đánh dấu rõ ràng, từng mục một được hoàn thành đúng lịch trình.

``Với cả cũng đến giờ nghỉ giải lao ngắn của chúng ta rồi,'' anh nói, gấp bảng ghi chú lại một cách gọn gàng.

``Nhưng mình còn chưa chạy thử cả bài hát đâu,'' tôi vội phản đối, cảm thấy chưa muốn dừng lại khi mọi thứ mới chỉ bắt đầu.

``Thế nên chúng ta cần phải nghỉ ngơi trước khi bắt đầu phần quan trọng đó,'' anh Kuon cười nhẹ, giải thích lý do. ``Cơ thể em cần thư giãn để chạy được toàn bộ bài hát một cách tốt nhất.''

Tôi với tay lấy chai nước trên bàn, lần này cầm đúng chiều. Ngay lúc đó, loa kiểm âm trên trần nhà phát ra một tiếng bíp ngắn, giọng đều đặn của Game vang lên: ``Hành vi cầm chai nước đã được cải thiện. Xác nhận thành công.''

``Cảm ơn ông rất nhiều về lời khen đó,'' tôi bĩu môi nói, cố gắng giữ vẻ mặt nghiêm trọng trong khi cầm chai nước lên, để mở nắp và uống một ngụm nhỏ.

\ct

Ngay khi chúng tôi vừa kết thúc buổi kiểm tra chuyển động và chuẩn bị bước vào phần tập chính của buổi tổng duyệt, cánh cửa phòng thu nhẹ nhàng mở ra, và Suzuno bước vào với một chiếc khay gỗ nhỏ được bọc khăn vải hoa trên tay. Trên khay xếp gọn gàng vài chiếc bánh mì kẹp nhân bơ mật ong được cắt nhỏ vừa ăn, những lát táo xanh tươi rói được xếp trong chiếc hộp nhựa trong, và một chiếc bình thủy tinh giữ nhiệt chứa trà xanh ấm mà em ấy thường pha cho cả nhà. 

``Em nghĩ là mọi người sẽ cần có gì đó để ăn uống trong lúc nghỉ giữa buổi tập nên đã chuẩn bị sẵn ạ,'' em ấy nói nhỏ, bước nhẹ nhàng qua khoảng trống trên sàn để không va vào các dây cáp nối từ máy tính ra màn hình lớn.

Anh Kuon nhanh chóng đứng dậy đón lấy chiếc khay từ tay em ấy, cẩn thận đặt nó lên chiếc bàn gỗ con sát cửa phòng, nơi có đủ không gian để mọi người có thể tự lấy đồ ăn khi cần. ``Cảm ơn em nhiều lắm, Suzuno. Bọn anh vừa mới định nghỉ giải lao đây.''

Tôi cũng hơi bất ngờ khi thấy em ấy xuất hiện ở phòng thu sớm thế, vốn nghĩ hôm nay em ấy có lịch học ở trường như mọi ngày. ``Suzuno-chan, hôm nay em không phải đi học sao? Chị tưởng em có tiết học buổi sáng mà?'' tôi hỏi, đặt cây bút chì đang cầm trên tay xuống cạnh danh sách kiểm tra.

``Hôm nay trường em nghỉ buổi sáng ạ. Em đã dậy sớm nấu ăn sáng và chuẩn bị bữa trưa cho mọi người trong nhà xong xuôi rồi mới mang đồ ăn nhẹ lên đây,'' con bé đáp, vừa nói vừa bước lại gần màn hình lớn ở cuối phòng, nơi mô hình 3D của tôi vẫn đang đứng yên lặng giữa sân khấu ảo được chiếu sáng. Đôi mắt em ấy tròn xoe nhìn ngắm những hiệu ứng ánh sáng lung linh trên sân khấu, rồi khẽ thốt lên: ``Trông đẹp quá chị ơi. Hôm nay chị định chạy thử toàn bộ bài hát phải không ạ?''

``Ừ, định thế. Đây chỉ là buổi tổng duyệt thôi, để kiểm tra xem mọi thứ có hoạt động ổn hết không trước buổi diễn chính,'' tôi vuốt lại mép tạp chí trên bàn cho thẳng thắn.

Suzuno quay lại nhìn tôi, ánh mắt hơi ngại ngần hỏi: ``Vậy em có thể ở lại ngồi xem một chút được không ạ? Em sẽ không gây trở ngại gì cho buổi tập đâu.''

Tôi quay sang nhìn anh Kuon, người chỉ nhún vai một cái, ý rằng để tôi tự quyết định việc em ấy có thể ở lại hay không. Tôi quay lại nhìn Suzuno, mỉm cười và gật đầu: ``Được chứ, cứ ở lại thoải mái. Nhưng nếu em thấy buồn chán thì cứ tự đi lên nhà chơi nhé, không cần phải ngồi qua hết buổi tập đâu.''

Em ấy vui vẻ ngồi xuống chiếc ghế đơn sát tường, xếp gọn gàng hai tay lên đầu gối như một học sinh ngoan trong lớp học, và đảm bảo rằng chiếc ghế của em ấy nằm hoàn toàn ngoài khu vực tập luyện của tôi trên sàn: ``Em nhất định sẽ không làm phiền đến chị và mọi người đâu ạ.''

Ngay sau đó, giọng của Game vang lên từ loa kiểm âm, báo cáo rằng ông ấy đã đánh dấu hoàn thành khung thời gian nghỉ giải lao mười phút trong lịch trình tổng duyệt của ngày hôm nay. Tôi đứng dậy đi lấy một chiếc bánh mì nhỏ từ khay, ăn nửa chiếc để lót dạ, trong khi Suzuno lấy chiếc bình giữ nhiệt rót trà ra hai cốc giấy, một cốc cho tôi và một cốc cho anh Kuon.

Hơi ấm từ cốc trà bốc lên nhẹ nhàng, mang theo mùi trà xanh thanh dịu quyện với vị ngọt của mật ong mà em ấy đã pha thêm: ``Em pha loại trà dịu thôi ạ, không pha quá đặc để không ảnh hưởng đến giọng hát của chị,'' Suzuno nói, đẩy cốc trà về phía tôi. ``Chị cứ uống từng ngụm nhỏ thôi nhé, đừng uống quá nhanh vì trà còn hơi nóng, dễ làm tổn thương cổ họng trước khi bắt đầu hát.''

Tôi nhận lấy cốc trà từ tay em ấy, lòng ấm áp vô cùng vì sự chu đáo của em: ``Cảm ơn em nhiều lắm, Suzu-chan. Em luôn nghĩ đến mọi thứ thật kỹ lưỡng.'' Tôi nâng cốc trà lên nhấp một ngụm nhỏ, cảm nhận vị ấm lan tỏa xuống cổ họng.

Anh Kuon lúc này đang ngồi xuống mở lại cuốn lịch trình đã xếp gọn trước đó, lướt qua các mục ghi chú bằng bút màu trên giấy. ``Sau khi hết thời gian nghỉ này, chúng ta sẽ tập trung chạy lại phần điệp khúc trước tiên. Nếu mọi thứ ổn thỏa rồi, chúng ta mới sẽ chạy thử toàn bộ bài hát từ đầu đến cuối,'' anh thông báo kế hoạch cho tôi, chỉ vào phần điệp khúc được tô màu vàng trong sơ đồ bài hát.

Tôi lặng lẽ nhìn vào trang giấy đó, lòng dấy lên một chút lo lắng. Điệp khúc chính là phần mà tôi đã gặp nhiều vấn đề nhất trong các buổi tập luyện trước đây, luôn bị lệch nhịp thở và bước chân mỗi khi đến đoạn đó, nhưng tôi đã không nói ra nỗi lo đó với ai cả. Ngay lúc đó, Suzuno khẽ đặt cốc trà của em ấy xuống bàn, như thể đọc được suy nghĩ trong lòng tôi, và nói nhỏ nhẹ: ``Chị đừng lo lắng quá, chị chắc chắn sẽ làm được mà.''

Tôi ngước nhìn em ấy, nở ra một nụ cười nhẹ để trấn an chính mình cũng như đáp lại sự động viên của em. ``Chị sẽ cố gắng hết sức mình.''

\ct

Bốn nhịp nhạc dạo đầu tiên vang lên, trầm ấm và rõ ràng từ hệ thống loa phòng thu, báo hiệu lượt thử nghiệm đầu tiên chính thức bắt đầu. Tôi đứng vững trên dấu chân trung tâm được sơn trắng mờ trên sàn, các cơ bắp vai đã thả lỏng hoàn toàn, mắt dán chặt vào điểm đánh dấu nhỏ màu xanh lá trên màn hình lớn, vị trí mà tôi đã tập trung nhìn vào trong suốt hàng chục buổi luyện tập trước đây. Từng câu hát, từng chuyển động trong đoạn mở đầu đều trôi đi hoàn toàn suôn sẻ, đúng y hệt như những gì tôi đã thuộc nằm lòng.

Tôi cất tiếng hát câu đầu tiên, giọng nói lan tỏa đều đặn trong không gian phòng thu. Trên màn hình, mô hình 3D của tôi cũng di chuyển cánh tay đồng bộ theo từng nhịp tay của tôi trong thế giới thực. Đến khi kết thúc câu cuối cùng của đoạn đầu, tôi thực hiện đúng động tác bước sang phải mà đã tập luyện hàng trăm lần, rồi quay trở lại vị trí trung tâm sân khấu đúng vào nhịp nhạc dự kiến. Ngay sau đó, giai đoạn chuyển tiếp của bài hát bắt đầu: tôi nâng tay lên cao, xoay toàn bộ thân người theo đúng sơ đồ, rồi tiến thêm một bước về phía trước. Mọi thứ dường như vẫn nằm trong tầm kiểm soát, cho đến khi điệp khúc chính tiến đến với những tiếng trống dồn dập hơn, tăng tốc độ nhịp điệu đột ngột.

Tốc độ các động tác của tôi cũng phải tăng theo để bắt kịp nhịp nhạc. Tôi bước chéo theo đúng đường đã đánh dấu, xoay vai mạnh mẽ và đưa cả hai tay lên cao như đã được thiết kế. Nhưng đúng ngay trước khi tôi mở miệng hát câu đầu tiên của điệp khúc, hơi thở của tôi đột ngột hụt đi; tôi không kịp nạp đủ không khí để phát ra âm thanh mạnh mẽ như cần thiết.

Kết quả là tôi vào câu hát chậm mất đúng nửa nhịp. Giọng hát tôi vẫn vang lên, nhưng nó không còn ăn khớp với những chuyển động của mô hình trên màn hình nữa. Cảm giác sai lệch đó quá rõ ràng, tôi buộc phải dừng lại giữa chừng.

Nhạc lập tức tắt ngay khi tôi dừng bước. ``Chúng ta thử lại từ đầu đoạn chuyển tiếp nhé,'' anh Kuon đề xuất một cách bình tĩnh từ phía bàn làm việc.

Tôi chỉ có thể gật đầu đồng ý, lòng hơi xao lạc. Suzuno không nói gì để an ủi hay phán xét, em ấy chỉ lặng lẽ đẩy cốc trà xanh ấm còn bốc hơi lại gần vị trí của tôi hơn một chút, như một hành động nhỏ bé ủng hộ. Tôi với tay lấy cốc, uống một ngụm nhỏ để làm ấm cổ họng, rồi quay sang hỏi chiếc loa kiểm âm: ``Có phải ông đã vào sai nhịp của bản nhạc rồi không, Game?''

``Không phải đâu,'' giọng của vị trợ lý ảo trả lời ngay lập tức. ``Câu hát của cô vào đúng hoàn toàn với nhịp trên bản phối gốc. Vấn đề nằm ở nhịp thở của cô đã thay đổi bất thường ngay trước đó.''

``Nghĩa là tôi đã lấy hơi quá muộn so với cần thiết?'' tôi hỏi lại, cố gắng phân tích tình huống.

``Dữ liệu từ cảm biến trên người cô cho thấy cô đã bắt đầu thực hiện động tác xoay người trước khi hoàn tất quá trình hít vào hơi thở,'' Game giải thích chi tiết.

Anh Kuon lúc này phóng to phần ghi chú trên máy tính bảng của mình, chỉ vào dòng chữ được tô sáng. ``Vấn đề không hề nằm ở các nốt nhạc hay câu hát đâu. Chính động tác xoay người đó đang chiếm mất hết khoảng thời gian quý báu mà em cần để lấy hơi cho câu hát tiếp theo.''

Tôi nhìn xuống bản sơ đồ chuyển động trải trên mặt bàn. Trên tờ giấy ấy, bước xoay người chỉ đơn giản là một dấu mũi tên mảnh, trông vô cùng đơn giản. Nhưng khi thực hiện nó ngoài thực tế, động tác đó buộc tôi phải thay đổi hướng toàn bộ thân thể, tập trung giữ thăng bằng trong suốt quá trình xoay, và chỉ sau đó mới có thể mở miệng hát\dots\ thực sự mất quá nhiều thời gian trong khi nhịp nhạc không chờ đợi ai.

``Chị có thể thử bỏ hoàn toàn động tác xoay đó đi không ạ?'' Suzuno nhẹ nhàng đưa ra một gợi ý từ phía ghế ngồi của em ấy.

Tôi lắc đầu ngay lập tức, không muốn thay đổi yếu tố quan trọng của vũ đạo. ``Đoạn chuyển tiếp đó nhất thiết phải có chuyển động để thu hút mắt người xem. Nhưng có lẽ\dots\ chắc chị có thể thay đổi cách thực hiện bước chân đó.''

Anh Kuon kéo chiếc ghế của mình lại gần bàn hơn, tạo không khí thảo luận tập trung. ``Chúng ta có thể ghi nhận ba phương án để thử nghiệm: giảm tốc độ khi thực hiện động tác xoay, thay đổi vị trí lấy hơi trong nhịp nhạc, hoặc hoàn toàn thay đổi cách bước chân để thực hiện động tác đó.''

Game ngay lập tức bổ sung: ``Tôi đề nghị chúng ta nên thử từng phương án một cách lần lượt, để có thể đo lường và so sánh kết quả một cách chính xác.''

Tôi một lần nữa gật đầu đồng ý với kế hoạch. Và như vậy, lượt thử thứ hai của chúng tôi chính thức bắt đầu.

\ct

Trong lần thử thứ hai này, tôi quyết định áp dụng phương án đầu tiên: giảm tốc độ khi thực hiện động tác xoay người. Nhưng kết quả thì không như mong đợi: toàn bộ động tác trông nặng nề, cứng nhắc hơn rất nhiều so với dự kiến ban đầu. Tôi có kịp lấy đủ hơi để hát câu đầu của điệp khúc, nhưng bước tiếp theo sau đó lại bị trễ mất nhịp.

Mô hình 3D trên màn hình vẫn theo sát từng chuyển động của tôi một cách chính xác, thế nhưng toàn bộ phần điệp khúc lại bắt đầu muộn hơn so với nhịp nhạc đang phát. Một lần nữa, tôi buộc phải dừng lại giữa chừng. ``Không ổn chút nào,'' tôi thốt lên, cảm giác hơi nản lòng.

``Không sao đâu,'' anh Kuon nhanh chóng an ủi. ``Ít nhất chúng ta cũng đã loại bỏ được một phương án không khả thi, đó cũng là một tiến bộ.''

Đến lượt thử thứ ba, tôi quyết định thử phương án thứ hai: cố gắng lấy hơi sớm hơn nhiều so với bình thường. Tôi hít vào một hơi sâu ngay trước khi bắt đầu thực hiện động tác xoay người. Nhưng sự thay đổi này lại tạo ra một vấn đề mới: hơi thở của tôi trở nên quá dài, bị dồn dập không cần thiết. Đến khi thực sự mở miệng hát, câu đầu tiên của điệp khúc nghe yếu hơn hẳn mức cần thiết. Tôi vô thức nhíu mày lại, không hài lòng với chất lượng giọng hát của mình.

``Chị có thể cố gắng giữ hơi đó cho đến hết câu hát được không?'' Suzuno lại một lần nữa đưa ra gợi ý, cố gắng tìm cách giải quyết vấn đề.

``Có thể giữ được,'' tôi đáp, ``nhưng như vậy thì các câu hát ở đoạn sau chắc chắn sẽ không còn đủ lực để thể hiện đúng cảm xúc của bài hát được.''

Ngay lúc đó, Game hiển thị lên màn hình phụ một biểu đồ âm lượng chi tiết, đi kèm với phân tích: ``Biên độ âm thanh của câu hát đã giảm mười hai phần trăm ở nửa sau của câu đầu tiên. Mức độ giảm này đủ lớn để khán giả có thể nhận ra.''

Tôi thở dài một cách mệt mỏi. ``Em không muốn đoạn cuối của bài hát bị yếu đi, mất hết sức mạnh vốn có.''

Anh Kuon khoanh tay trước ngực, chìm vào suy nghĩ trong vài giây, rồi mới đề xuất phương án cuối cùng: ``Vậy thì chúng ta thử thay đổi hoàn toàn động tác bước chân. Nếu em thực hiện động tác xoay bằng các bước chân nhỏ hơn thay vì một bước dài duy nhất, phần thân trên của em sẽ ổn định hơn nhiều, và em có thể lấy hơi một cách tự nhiên mà không bị gián đoạn.''

Tôi nhìn xuống các dấu chân được vẽ trên sàn phòng thu. Bước xoay mà tôi đang sử dụng hiện tại được đánh dấu bằng một đường chéo dài, kéo từ vị trí này sang vị trí khác. Tôi thử nhẩm lại trong đầu: nếu chia nó thành hai bước ngắn riêng biệt - một bước đơn giản sang bên, rồi một bước quay trở lại vị trí trung tâm - tôi sẽ không cần phải nghiêng người quá nhiều để giữ thăng bằng.

Suzuno dường như đọc được suy nghĩ của tôi, em ấy đứng dậy từ ghế và thực hiện chậm rãi động tác đó để tôi có thể quan sát rõ từng bước. ``Chị có thể thử xoay vai trước tiên, rồi mới thực hiện chuyển động của chân ạ,'' em ấy hướng dẫn một cách nhẹ nhàng.

Tôi làm theo đúng như em ấy chỉ. Lần đầu tiên thực hiện, toàn bộ động tác vẫn còn khá cứng nhắc, chưa tự nhiên. Nhưng đến lần thứ hai, trọng tâm cơ thể của tôi đã ổn định hơn nhiều. Và ở lần thứ ba luyện tập động tác mới, tôi đã có thể lấy hơi một cách hoàn toàn tự nhiên mà không cần phải suy nghĩ hay tập trung quá nhiều vào việc đó.

``Bây giờ chúng ta thử thực hiện lại với bản nhạc nhé,'' anh Kuon đề xuất, nhìn thấy sự tiến bộ rõ rệt.

\ct

Nhạc bản được bật lại từ điểm đầu đoạn chuyển tiếp, đúng như chúng ta đã thống nhất. Tôi xoay vai theo đúng hướng đã tập, bước chân nhẹ nhàng dịch sang bên phải rồi quay lại vị trí trung tâm sân khấu. Động tác mới ngắn gọn hơn nhiều so với vòng xoay lớn ban đầu, nhưng vẫn giữ được sự mềm mại, uyển chuyển của vũ đạo, không hề bị cứng nhắc hay đơn điệu.

Tôi hít vào một hơi sâu ngay trong khoảng lặng ngắn sau tiếng trống đầu tiên của đoạn chuyển tiếp. Khi điệp khúc chính thức bắt đầu, tôi mở miệng hát đúng nhịp, không bị chậm lại nửa giây nào. Lần này mọi thứ đều diễn ra đúng như mong đợi.

Trên màn hình lớn, mô hình 3D của tôi di chuyển đồng bộ hoàn hảo với từng cử chỉ của tôi ngoài đời thực. Tôi hát trọn vẹn câu đầu tiên của điệp khúc, và lượng hơi thở còn dư vẫn đủ để tôi tiếp tục câu hát tiếp theo mà không bị hụt hơi.

Khi đến cuối đoạn điệp khúc, tôi giữ được nốt ngân dài một cách tự nhiên, không cần phải gồng cổ hay cố gắng quá sức. Nhạc lập tức tắt đi sau khi nốt cuối cùng tan biến. Tôi đứng yên tại chỗ, mắt dán chặt vào màn hình, chờ đợi thông báo kết quả kiểm tra từ Game.

``Độ lệch giữa giọng hát và chuyển động đã giảm xuống dưới ngưỡng đo lường tối thiểu,'' giọng đều đặn của Game vang lên từ loa, báo cáo kết quả một cách chính xác.

Anh Kuon gật đầu hài lòng từ phía bàn làm việc. ``Tốt hơn rất nhiều so với mấy lần trước rồi đấy.''

Suzuno vỗ tay nhẹ nhàng, nụ cười rạng rỡ nở trên môi. Tôi thở một hơi dài ra, lúc này mới nhận ra hai vai mình đã căng cứng suốt thời gian thực hiện động tác, như thể đang nắm chặt một thứ gì đó quý giá.

``Em làm được rồi, Delu-san!'' giọng em ấy đầy vui mừng.

``Chỉ mới qua được đoạn này thôi,'' tôi đáp, vẫn còn cảm thấy chưa hoàn toàn hài lòng.

``Đoạn này cũng là một phần quan trọng của bài hát mà, đã vượt qua được thì tốt lắm ạ,'' Suzuno nhanh chóng động viên.

Tôi không biết phải đáp lại lời khen của em ấy như thế nào, nên chỉ có thể mỉm cười. Ngay sau đó, Game hiển thị bảng kết quả chi tiết lên màn hình phụ, với các chỉ số được đánh dấu rõ ràng.

``Tôi đề nghị mọi người tạm nghỉ ba phút trước khi bắt đầu chạy thử toàn bộ bài hát từ đầu đến cuối,'' vị trợ lý ảo đề xuất.

Anh Kuon quay sang nhìn tôi, ý nhắn nhủ rõ ràng. ``Lần này nghe lời Game đi, đừng cố gắng quá sức.''

``Em biết rồi,'' tôi gật đầu đồng ý, hiểu được sự lo lắng của anh.

Tôi bước về phía ghế của mình, ngồi xuống rồi nhấp một ngụm nhỏ trà xanh ấm từ cốc giấy. Hơi ấm lan tỏa xuống cổ họng, giúp các cơ bắp căng thẳng từ từ thả lỏng.

\ct

Đến giữa trưa, Kson ghé qua phòng thu sau khi vừa hoàn thành xong công việc ở tầng dưới. Cô bước vào đúng lúc tôi đang ngồi xem lại video của buổi tổng duyệt vừa rồi, phân tích từng khung hình để tìm ra những điểm chưa ổn. 

``Sao rồi bé Delu? Có gì khó khăn không?'' cô hỏi, tiến lại gần bàn làm việc.

``Đoạn điệp khúc ban đầu hơi vất vả một chút. Bọn em vừa mới chỉnh sửa bước xoay người, thay vì một vòng xoay lớn thì chia thành hai bước nhỏ hơn,'' tôi giải thích, chỉ vào màn hình đang dừng ở khoảnh khắc tôi thực hiện động tác mới.

Kson đứng sau lưng ghế của tôi, mắt dán vào màn hình. ``Cho chị xem lại đoạn đó được không? Tò mò muốn xem động tác mới thế nào.''

Anh Kuon nhấn nút phát lại từ đầu đoạn chuyển tiếp. Kson xem im lặng suốt đoạn nhạc, mắt theo dõi sát sao từng chuyển động của mô hình 3D trên sân khấu ảo. Khi đoạn nhạc kết thúc, cô nghiêng đầu một chút, như thể đang đánh giá.

``Động tác mới trông tự nhiên hơn hẳn, không bị gượng ép như vòng xoay cũ,'' chị nhận xét.

``Em cũng sợ rằng nó quá đơn giản, không đủ thu hút sự chú ý của khán giả,'' tôi thừa nhận nỗi lo của mình.

``Sân khấu đâu cần phải lấp đầy mọi giây bằng những động tác lớn lao đâu. Có lúc một bước nhỏ đúng nhịp còn rõ ràng, còn gây ấn tượng hơn cả một vòng xoay phức tạp đó!'' Kson nói, chỉ vào màn hình. ``Đừng nghĩ đơn giản là tầm thường.''

Tôi nhìn lại hình ảnh trên màn hình, nơi tôi vừa quay về vị trí trung tâm sân khấu. Kson chỉ vào đúng khoảnh khắc đó. ``Đoạn này khán giả sẽ nhìn thấy mặt em rõ nhất, mọi ánh mắt đều đổ dồn về đây. Em không cần phải cố gắng lấp đầy mọi khoảng trống bằng chuyển động, để mọi thứ có khoảng nghỉ.''

``Chị nghĩ em nên giữ khoảng lặng ở đây à? Để nó tự nhiên thôi?'' tôi hỏi lại, muốn chắc chắn hiểu ý của cô.

``Ừ, để câu hát cũng có chỗ để thở. Không phải lúc nào cũng phải chuyển động liên tục mới là hay,'' Kson nhấn mạnh.

Tôi nhìn cô, ngạc nhiên trước những phân tích sắc sảo. ``Chị nói nghe như một chuyên gia sân khấu thực thụ ấy.''

``Chị có nhiều năm kinh nghiệm đứng trước máy quay, đứng trước đám đông mà. Từng này đã là gì đâu,'' Kson cười nhún vai.

``Em còn tưởng chị đến đây để trêu em như mọi khi chứ,'' tôi nói đùa.

``Chị có thể làm cả hai mà, vừa góp ý chân thành vừa trêu em một chút,'' Kson cười lớn, rồi vừa ngồi xuống chiếc ghế trống cạnh bàn.

Anh Kuon khoanh tay nhìn cô, nhíu mày. ``Kson, em có thể xem lịch tổng duyệt mà không chiếm luôn ghế của Delu không? Và đừng ngồi lên đó nữa.''

Kson nhìn xuống chiếc ghế dưới mình, rồi nhìn xuống mớ giấy tờ mà mình vừa ngồi lên, rồi nhích sang bên một chút. ``Chị chỉ đang hỗ trợ tinh thần cho em ấy thôi mà.''

``Hỗ trợ tinh thần đâu cần phải ngồi lên cả bảng ghi chú của người ta,'' anh Kuon lắc đầu bất lực.

Kson vội vàng đứng dậy khỏi xấp giấy, một số tờ đã bị gấp nhẹ cạnh góc. Ngay lúc đó, Game lên tiếng: ``Một số ghi chú trên giấy đã bị gấp lại, nhưng tôi sẽ khôi phục bản điện tử để mọi người vẫn sử dụng được.''

Tôi cười thành tiếng, cảm thấy không khí trong phòng từ từ trở nên nhẹ nhàng, thoải mái hơn sau những căng thẳng của buổi tập.

\ct

Khi khoảng nghỉ kết thúc, chúng tôi chuyển sang giai đoạn cuối cùng của buổi tổng duyệt: chạy thử toàn bộ bài hát từ đầu đến cuối, không dừng lại giữa chừng. Tôi bước chậm rãi đến vị trí trung tâm của khu vực tập luyện, đôi chân đặt trùng khớp với dấu chân sơn trắng mờ trên sàn gỗ, cảm nhận sự quen thuộc của không gian này. Trên màn hình lớn, mô hình 3D của tôi cũng đứng yên trong tư thế tương tự, chờ đợi tín hiệu bắt đầu. Loa kiểm âm trên trần nhà bật lên, giọng đều đặn của Game bắt đầu phát ra chuỗi đếm ngược đã quá quen thuộc.

Ba.

Hơi thở của tôi chậm lại, các cơ bắp vai thả lỏng dần để chuẩn bị cho các chuyển động sắp tới.

Hai.

Tôi nâng nhẹ cằm, mắt dán chặt vào điểm đánh dấu xanh lá nhỏ trên màn hình.

Một.

Nhạc lập tức vang lên, bốn nhịp dạo đầu trầm ấm lan tỏa khắp phòng thu, đánh thức mọi giác quan của tôi. Đoạn mở đầu trôi qua một cách êm thấm, mọi chuyển động đều đi đúng theo nhịp mà tôi đã thuộc nằm lòng. Tôi giữ vững ánh mắt vào điểm nhìn đã chọn, không để ánh mắt đảo lộn làm mất thăng bằng.

Tay trái nâng lên đúng vào nhịp thứ tám của đoạn mở đầu, cử chỉ mềm mại như đã luyện tập hàng trăm lần. Đến câu hát thứ hai, tay phải của tôi chạm nhẹ vào vùng ngực, một cử chỉ nhỏ để truyền tải cảm xúc của lời bài hát. Trên màn hình, mô hình 3D thực hiện đồng bộ từng chuyển động, không có độ trễ nào cảm nhận được, mọi thứ dường như đang diễn ra hoàn hảo.

Khi giai đoạn chuyển tiếp đến, tôi tự động nhớ ra bài học từ buổi sáng: phải lấy hơi sớm hơn. Tôi hít một hơi sâu ngay sau tiếng trống đầu tiên của đoạn chuyển tiếp, xoay vai theo đúng hướng đã tập luyện, bước chân ngắn sang bên phải rồi quay lại vị trí trung tâm trước khi điệp khúc chính thức bùng nổ.

Tôi mở miệng hát câu đầu tiên của điệp khúc mà không hề gặp tình trạng hụt hơi như trước đây. Mọi thứ diễn ra đúng như mong đợi, đến câu thứ hai tôi còn kéo dài nốt nhạc lâu hơn một chút so với dự tính ban đầu, và vẫn giữ được nhịp độ ổn định.

Nhưng đến đoạn nghỉ ngắn giữa bài, một sai sót nhỏ xảy ra: tôi đưa tay lên chậm hơn đúng nửa nhịp. Ngay lập tức, tôi nhận ra vấn đề khi nhìn vào màn hình: mô hình 3D đã bắt đầu thực hiện chuyển động nâng tay trong khi bàn tay của tôi ở ngoài đời thực vẫn còn đang ở vị trí thấp. Cảm giác sai lệch ấy hiện rõ trước mắt, và phản xạ đầu tiên của tôi là dừng lại, yêu cầu bắt đầu lại như mọi lần trước. Nhưng tôi vẫn còn câu hát tiếp theo cần hoàn thành, một tiếng nói nhỏ trong đầu nhắc nhở. Tôi giữ vững giọng hát, nhanh chóng đưa tay lên đúng vào nhịp sau, kịp thời để mô hình bắt kịp lại chuyển động của tôi.

Tôi tiếp tục vượt qua đoạn cuối của bài hát, giữ được nốt ngân cuối cùng một cách trọn vẹn, tự nhiên. Khi nốt nhạc cuối cùng tan biến, loa tự động tắt âm thanh. Cả căn phòng rơi vào im lặng tức thì, không ai lên tiếng ngay cả khi tôi vẫn đang đứng ở giữa phòng, cố gắng điều chỉnh lại hơi thở của mình sau khi hoàn thành toàn bộ bài hát.

Tôi nhìn sang phía bàn làm việc, thấy anh Kuon đang chăm chú xem bảng dữ liệu trên máy tính bảng, đôi mắt phân tích từng chỉ số. Suzuno ngồi ở ghế sát tường, ánh mắt trông chờ nhìn về phía tôi, còn Kson thì khoanh tay dựa vào tường, nụ cười trên môi khó đoán được ý nghĩa.

Game là người đầu tiên phá vỡ bầu không khí im lặng, giọng đều đặn của ông phát ra từ loa: ``Toàn bộ bài hát đã được hoàn thành. Có một tín hiệu chuyển động của tay bị lệch nửa nhịp tại đoạn nghỉ, nhưng cô đã tự điều chỉnh mà không làm gián đoạn phần trình diễn ca hát.''

Tôi quay lại nhìn màn hình, nơi Game đã tua lại đoạn video ghi lại lượt vừa rồi. Trong hình ảnh, mô hình thực sự có một chút chậm trễ ở động tác nâng tay, nhưng toàn bộ phần trình diễn vẫn tiếp tục diễn ra liền mạch, không bị phá vỡ.

``Em có thể sửa lỗi đó ở những lần tập sau,'' anh Tổng nói lên tiếng, gật đầu hài lòng. ``Nhưng lượt trình diễn này đã cho thấy em biết cách xử lý tình huống khi có sai lệch xảy ra, đó là một kinh nghiệm quý giá. Tốt lắm.''

Tôi chậm rãi gật đầu, lòng vẫn còn cảm thấy ngạc nhiên về phản ứng của bản thân. ``Trước đây em cứ nghĩ rằng nếu có bất cứ điều gì sai sót, thì phải dừng lại và làm lại từ đầu, không có lựa chọn nào khác.''

Kson nhún vai, ánh mắt có vẻ đồng tình. ``Trên sân khấu thực tế, đôi khi cách sửa lỗi tốt nhất chính là tiếp tục đi tiếp, không được phép dừng lại giữa chừng. Em đã làm tốt điều đó.''

Lúc này, Suzuno đứng dậy từ ghế của mình, bước lại gần và đưa cho tôi chiếc khăn nhỏ mà em ấy đã chuẩn bị. ``Chị đã hoàn thành hết toàn bộ bài hát rồi ạ.''

Tôi nhận lấy chiếc khăn từ tay em, lau nhẹ những giọt mồ hôi trên trán, cảm giác thành tựu dâng lên trong lòng. ``Ừ. Cuối cùng chị cũng đã hoàn thành được.''

\ct

Chúng tôi quyết định thực hiện thêm một lượt tổng duyệt nữa để củng cố sự ổn định. Lần này, tôi không còn cố ép từng động tác phải giống hệt từng khung hình trong video hướng dẫn như trước đây. Tôi vẫn giữ nguyên những dấu mốc quan trọng, những bước chân và nhịp hát không thể thay đổi, nhưng để cho vai và cánh tay của mình chuyển động một cách tự nhiên hơn, không bị gò bó bởi khuôn mẫu.

Tôi vẫn bước đến đúng các vị trí đã được đánh dấu trên sàn. Vẫn vào câu hát đầu tiên của điệp khúc đúng nhịp. Vẫn giữ được nốt ngân cuối cùng một cách trọn vẹn.

Khi đến đoạn nâng tay mà lượt trước tôi đã bị chậm nửa nhịp, tình huống tương tự lại xảy ra. Nhưng lần này tôi không hề có ý định dừng lại. Tôi tiếp tục giữ vững giọng hát, rồi đưa tay lên đúng vào nhịp kế tiếp, như thể đó là một phần của kế hoạch từ đầu.

Lần này, chuyển động đó trông hoàn toàn giống một cử chỉ có chủ ý, một sáng tạo nhỏ trong quá trình trình diễn chứ không phải là một sai sót. Khi bản nhạc kết thúc, tôi nghe tiếng vỗ tay nhẹ nhàng của Suzuno từ phía ghế ngồi. Kson cũng vỗ tay hai lần, tiếng vỗ tay khá lớn thể hiện sự công nhận.

Anh Kuon đánh dấu các kết quả vào bảng lịch trình, rồi ngước nhìn lên. ``Lượt thứ hai này ổn định hơn nhiều so với lượt đầu. Vẫn còn một sai lệch nhỏ ở động tác tay, nhưng nó không hề ảnh hưởng đến chất lượng tổng thể của buổi trình diễn.''

Tôi nhìn lại sân khấu ảo trên màn hình lớn, nơi các hiệu ứng ánh đèn đang tắt dần, để lại mô hình 3D của tôi đứng yên giữa nền tối. Một cảm giác lạ lùng xuất hiện trong lòng: tôi không hề thấy mình cần phải chạy lại ngay lập tức để sửa chữa mọi thứ. Cảm giác này thậm chí còn làm tôi ngạc nhiên hơn cả việc vừa hoàn thành được toàn bộ bài hát.

``Chị có mệt không ạ?'' Suzuno bước lại gần, hỏi nhỏ.

``Có một chút,'' tôi thừa nhận, các cơ bắp trên người thực sự đang cảm thấy mệt mỏi sau một buổi tập tập trung.

``Vậy thì chúng ta nghỉ ngơi thôi nhé,'' em ấy đề xuất.

Tôi quay đầu nhìn chiếc đồng hồ treo trên tường. Kim đồng hồ đã chỉ gần ba giờ chiều, buổi tập kết thúc sớm hơn nhiều so với kế hoạch ban đầu mà chúng tôi đã vạch ra.

Kuon đóng cuốn bảng lịch lại một cách gọn gàng. ``Chúng ta đã thu thập đủ tất cả dữ liệu cần thiết để thực hiện các chỉnh sửa lần cuối vào ngày mai rồi.''

``Em có cần phải chạy thêm một lượt nào nữa trong hôm nay không ạ?'' tôi hỏi, vẫn còn hơi ngạc nhiên về việc kết thúc sớm này.

``Không cần đâu. Hôm nay em đã có một lượt trình diễn hoàn chỉnh, đủ để chúng ta làm việc rồi,'' anh Kuon trả lời chắc chắn.

Game ngay lập tức xác nhận thêm ý kiến của anh ấy: ``Việc lặp lại thêm các lượt tập khi cơ thể đang trong tình trạng mệt mỏi có thể làm giảm độ chính xác của các chuyển động, không mang lại lợi ích gì.''

Tôi bật cười trước sự đồng lòng của cả hai người: ``Được rồi. Vậy hôm nay chúng ta sẽ nghỉ tại đây thôi ạ.''

\ct

Khi Kson và Game bắt đầu dọn dẹp đống bảng ghi chú, máy tính bảng và các dây cáp loằng ngoằng trên bàn, tôi vẫn đứng yên một chỗ, chẳng muốn rời đi ngay. Ánh mắt tôi dán chặt vào màn hình lớn phía trước, nơi sân khấu ảo giờ đây đã tắt hết đèn nền, chỉ còn lại mô hình 3D của tôi đứng im lìm, hai tay buông thõng tự nhiên hai bên thân người, như thể nó cũng đang nghỉ ngơi sau buổi tập căng thẳng kéo dài cả ngày.

Tôi không thể ngừng hồi tưởng lại lượt chạy cuối cùng vừa rồi. Không phải là những nốt nhạc cao thấp mà tôi đã hát đúng hay sai, không phải là số lần bước chân đã đặt trùng khớp với các dấu vạch sơn trên sàn. Điều níu giữ suy nghĩ của tôi chính là khoảnh khắc tôi đã quyết định không dừng lại sau khi lỡ nhịp nâng tay. Cái phản xạ đầu tiên muốn hét ra lời ``dừng lại'' đã bị tôi kìm nén, và thay vào đó là tiếp tục hát, tiếp tục di chuyển, sửa lỗi một cách tự nhiên mà không làm gián đoạn cả buổi trình diễn. Đó là điều tôi chưa bao giờ làm được trước đây.

Tôi rời mắt khỏi màn hình, lục cặp túi đeo chéo lấy ra cuốn sổ nhỏ mà tôi luôn mang theo. Mở trang cuối cùng, nơi còn trống một khoảng nhỏ, tôi cầm bút mực, viết thật chậm một dòng ngắn gọn: ``Nếu lỡ sai, cứ hát tiếp.'' Bút mực chảy đều, những chữ viết hiện lên rõ ràng, như thể tôi đang khắc vào đó bài học quý giá mà mình mới học được.

Đột nhiên, có một hơi thở ấm áp len lỏi qua vai tôi, giọng Kson vang lên nhẹ nhàng từ phía sau: ``Câu đó hay đấy. Nên dán nó cạnh màn hình trong phòng thu lúc nào cũng thấy.''

Tôi quay cổ nhìn chị, cười nhẹ: ``Em chỉ viết để nhắc mình thôi, sẽ nhớ được mà.''

``Người ta mới hay viết ra giấy vì sợ mình quên đi những điều quan trọng đấy,'' Kson nhún vai, bước qua ngồi xuống chiếc ghế cạnh bàn, tay vỗ nhẹ lên gỗ mặt bàn. ``Đừng tự tin quá vào trí nhớ của mình, khi lên đài căng thẳng cái gì cũng có thể quên sạch.''

Tôi nhìn theo bóng cô, sau đó lật lại trang sổ, gật đầu đồng ý: ``Vậy thì em sẽ in ra dán cạnh màn hình như chị gợi ý.''

Ở phía đầu kia của căn phòng, Suzuno đang nhẹ nhàng dọn dẹp chiếc khay nhựa đựng các cốc trà và bánh ăn dặm suốt buổi. Chỉ còn lại chiếc bình thủy tinh đựng trà xanh, đã vơi đi gần một nửa, nằm đơn độc trên góc khay. Em ấy vừa xếp lại các chiếc khăn giấy bẩn, vừa quay sang nói với tôi: ``Em để phần bánh còn lại vào tủ lạnh trong bếp nhé, chị muốn ăn thì mai lấy.''

``Cảm ơn em nhiều lắm, Suzuno-chan,'' tôi đáp lại, lòng ấm áp bởi sự chu đáo của em ấy.

Lúc này, Kuon đã bật chế độ ngủ cho chiếc bảng ghi hình, màn hình tối dần đi. Anh nhìn sang tôi, giọng nói thoải mái hơn nhiều so với suốt buổi tập: ``Ngày mai chúng ta chỉ cần chỉnh lại tín hiệu theo dõi chuyển động của cánh tay cho mượt mà hơn, rồi chạy lại đoạn cuối bài hát để kiểm tra lần cuối là xong. ''

``Thật vậy sao? Nghe nhẹ nhàng hơn hẳn những gì em tưởng,'' tôi thở phào, cảm giác một gánh nặng nhỏ vừa rơi khỏi vai.

``Đương nhiên rồi, vì hôm nay em đã giải quyết xong phần khó nhằn nhất rồi mà,'' Kuon cười, vỗ tay vào đống bảng ghi chú đã xếp gọn gàng. ``Những gì còn lại chỉ là việc hoàn thiện nhỏ thôi.''

Trước khi bước ra khỏi phòng thu, tôi quay người lại, hướng mắt về màn hình lần cuối. Sân khấu ảo vẫn ở đó, im lìm nhưng tràn đầy sức sống, sẵn sàng chờ đợi chúng tôi quay lại vào ngày mai cho buổi tập cuối cùng. Tôi với tay ra, nhấn công tắc tắt hết các bóng đèn trong phòng, để căn phòng rơi vào bóng tối nhẹ nhàng.

\ct

Tối hôm ấy, sau khi tắm rửa sạch sẽ và thay bộ quần áo ở nhà thoải mái, tôi mang theo cuốn sổ nhỏ của mình đi xuống tầng trệt, nơi phòng khách chung của căn nhà luôn ấm áp. Bên trong, Ryuuhou đang ngồi trên ghế dài, tay vừa thay dây đàn guitar mới cho chiếc đàn acoustic của anh ấy, đôi tay khéo léo luồn dây qua các cầu đàn. Ở chiếc bàn trà đối diện, Kaigou đang chăm chú xem một bản báo cáo dài trên máy tính bảng, thỉnh thoảng lại chạm nhẹ màn hình để lướt sang trang tiếp theo.

Giữa tấm thảm lớn trải giữa phòng, Naomi đang ngồi xếp bằng, cầm bút màu vẽ say sưa trên một tờ giấy khổ lớn. Em ấy đang vẽ một sân khấu nhỏ xinh, với những chiếc đèn tròn đủ màu sắc vẽ quanh viền, trông thật sống động. 

Suzuno ngồi cạnh Naomi, vừa đổ trà nóng vào các tách trên bàn, vừa ngẩng lên nhìn tôi bước xuống. Em ấy mỉm cười, đặt chiếc tách trà vừa đổ vào vị trí trống dành cho tôi: ``Buổi tập hôm nay của chị thế nào rồi ạ? Có mệt không?''

``Tốt hơn rất nhiều so với những gì em đã mong đợi,'' tôi ngồi xuống cạnh em ấy, nhấp một ngụm trà ấm làm dịu cổ họng. ``Không ngờ lại hoàn thành sớm được thế.''

``Chị đã hát hết cả bài chưa ạ?'' Naomi ngẩng đầu lên từ bức tranh, đôi mắt trong sáng hỏi tôi.

``Ừ, chị hát được hai lượt hoàn chỉnh cả bài rồi,'' tôi gật đầu, mỉm cười với em nhỏ.

Con bé vui mừng cười toe toét, giơ bức tranh của mình lên cao, chỉ vào một góc trang giấy: ``Naomi cũng vẽ khán giả cho chị đây này. Tất cả mọi người đều đang vỗ tay cho chị đó!''

Tôi cúi người xuống, nhìn kỹ bức tranh đầy màu sắc. Những hình người nhỏ xíu với chiếc đầu tròn được vẽ đứng thành hàng dài trước sân khấu, mỗi người lại cầm một chiếc đèn hiệu màu khác nhau từ xanh dương, đỏ, vàng, tím\dots\ trông như một buổi hòa nhạc thực thụ.

Tim tôi bỗng dưng ấm áp vô cùng, tôi chỉ vào bức tranh, hỏi nhỏ em bé: ``Chị có thể giữ bức vẽ này của Naomi không? Chị muốn dán nó vào phòng của chị.''

Naomi gật đầu lia lịa, má ửng hồng vì vui: ``Dạ được ạ! Chị cứ lấy đi nhé.''

Tôi cẩn thận đặt bức tranh nhỏ cạnh cuốn sổ ghi chú của mình trên bàn. Lúc này, Kson vừa bước ra từ bếp, tay cầm lấy chiếc bánh macaron còn sót lại trên đĩa, cắn một nửa rồi ngước nhìn tôi: ``Ngày mai còn phải quay lại phòng thu tổng duyệt tiếp à?''

``Ừ, chỉ còn vài việc nhỏ thôi,'' tôi đáp lại.

``Vậy thì tối nay ăn uống xong xuôi là đi ngủ sớm đi, đừng thức khuya xem phim nữa,'' chị nhắc nhở, vừa nhai bánh vừa ngồi xuống bậu cửa sổ.

Tôi cười trêu cô: ``Chị đang nói với em hay là nhắc nhở chính bản thân chị đấy? Hôm qua ai thức đến hai giờ sáng xem bộ phim tình cảm kia?''

Kson bị tôi nói trúng tim đen, chỉ có thể cười khúc khích rồi gật đầu thừa nhận: ``Rồi, biết rồi. Mới tí tuổi đầu mà lẻo mép ghê thế.''

Tôi bật cười thành tiếng, cảm nhận không khí ấm áp tràn đầy trong căn phòng. Không ai vội vàng hỏi tôi đã hát bài gì, hay buổi trình diễn hôm nay có hoàn hảo đến đâu. Họ chỉ lắng nghe khi tôi muốn chia sẻ, và tối nay, tôi đã kể hết mọi chuyện về cái nửa nhịp tay bị chậm ở đoạn nghỉ, về cách bước chân mới mà chúng tôi đã thay đổi, và quan trọng nhất, là việc tôi đã không dừng lại khi có lỗi xảy ra.

Kaigou ngẩng lên từ chiếc máy tính bảng, gật đầu đồng tình với những gì tôi kể: ``Lần sau em sẽ nhớ được bài học đó thôi.''

``Ai biết được đâu, có thể đến lần sau em vẫn sẽ lỡ nhịp, vẫn sẽ mắc sai lầm,'' tôi cười, thừa nhận thực tế.

``Thế thì em lại tiếp tục đi tiếp thôi, có sao đâu,'' Kaigou nhún vai, nói một cách đơn giản mà chứa đựng cả một sự thấu hiểu.

Tôi nhìn xuống cuốn sổ trên bàn, rồi lại nhìn những người bạn xung quanh, và chỉ có thể nói một câu: ``Đúng rồi, vậy thì cứ tiếp tục đi thôi.''

\ct

Ngày hôm sau, tôi bước vào phòng thu với một tâm trạng hoàn toàn khác. Tay tôi cầm chiếc khăn tắm mới, và đôi giày tập luyện đã được buộc dây cẩn thận, chắc chắn. Trước khi bắt đầu, tôi lấy tờ giấy đã in dòng chữ hôm qua, dán nó ngay cạnh mép dưới màn hình lớn, nơi mà mỗi khi đứng ở giữa sân khấu, tôi đều có thể nhìn thấy: ``Nếu lỡ sai, cứ hát tiếp.'' Dòng chữ đen trên nền giấy trắng đơn giản, nhưng lại mang đến cho tôi một sự bình yên khó tả.

Tôi bước đến vị trí trung tâm của khu vực tập luyện, đôi chân đặt đúng vào các dấu chân sơn trắng mờ trên sàn. Mắt tôi nhìn vào mô hình 3D của mình trên màn hình, nó cũng đang đứng yên chờ đợi. Giọng nói đều đặn của Game vang lên từ loa trần, bắt đầu chuỗi đếm ngược quen thuộc: ``Ba.''

Hơi thở của tôi chậm lại, các cơ bắp trên vai từ từ thả lỏng, không còn sự căng cứng như những buổi tập trước. ``Hai.''

Tôi nâng nhẹ cằm, mắt nhìn vào điểm đánh dấu quen thuộc trên màn hình, tâm trí hoàn toàn tập trung. ``Một.''

Nhạc lập tức vang lên, những nốt đầu tiên lan tỏa khắp căn phòng. Tôi hít vào một hơi sâu đúng vào nhịp nghỉ của bản nhạc, chân tôi bước nhẹ nhàng dịch chuyển sang bên phải theo đúng động tác đã tập luyện. Khi đến đoạn điệp khúc chính, tôi cất tiếng hát một cách tự nhiên, không còn sự gắng sức hay lo lắng nào.

Lần này, tôi không hề nghĩ đến việc phải hoàn hảo trong từng chi tiết nhỏ. Tôi không ép bản thân phải sao chép đúng từng khung hình trong video hướng dẫn. Tôi chỉ nghĩ đến câu hát tiếp theo cần truyền tải cảm xúc thế nào, rồi lại đến câu kế tiếp, để âm nhạc dẫn dắt cơ thể mình di chuyển.

Khi nốt nhạc cuối cùng của bài hát tan biến, tôi vẫn đứng yên giữa căn phòng, tai nghe âm thanh cuối cùng vơi dần trong tai nghe. Không có tiếng vỗ tay vang dội của hàng ngàn khán giả trên sân khấu thực tế. Chỉ có giọng Game đọc ra các chỉ số kết quả, anh Kuon ngồi ở bàn làm việc ghi chú lại những điểm cần chỉnh sửa nhỏ, Suzuno đứng ở cửa phòng mỉm cười nhẹ nhàng với tôi, và Kson giơ ngón cái lên như một lời chúc mừng.

Đối với tôi, vậy là đã quá đủ cho buổi tổng duyệt hôm nay. Những gì còn lại, tất cả những cảm xúc, sự chuẩn bị cuối cùng, tôi sẽ giữ gìn và dành trọn nó cho ngày trình diễn chính thức, khi bước chân lên sân khấu lớn trước ánh mắt của mọi người.

%==========================================TRUYỆN PHỤ=========================================

\omk

Sau buổi tổng duyệt cuối cùng, khi mọi người đã dọn dẹp xong các thiết bị và chuẩn bị rời phòng thu, Kson bỗng dừng lại trước màn hình lớn vừa tắt. Chị quay sang nhìn tôi, người đang xếp dây tai nghe vào túi, rồi đột nhiên bước ra giữa khoảng trống sàn nhà nơi tôivừa đứng tập luyện, hỏi: ``Này Delu-chan, cái bước xoay vai rồi bước nhỏ em làm hôm qua đúng không? Chị thấy nó đơn giản mà hiệu quả quá, để chị thử xem.''

Vừa nói, Kson vừa nghiêng vai theo đúng hướng mà tôi đã tập, chân bước ngắn sang bên phải như đã xem hàng chục lần trong buổi tập. Nhưng rồi vì tính toán sai khoảng cách, gót chân chạm phải chân ghế sắt phía sau, suýt nữa thì ngửa người ra sau.

Anh Kuon, người đang đóng hộp các máy tính bảng, nhanh tay kéo chiếc ghế ra khỏi lối đi, lắc đầu bất lực. ``Này cẩn thận đi. Sân khấu chính của em rộng gấp mười lần cái phòng này, chứ ở đây mà còn va vào ghế thì lên đó làm sao?''

Kson vỗ tay vào vai, chỉnh lại tư thế đứng thẳng, mặt chẳng hề xấu hổ mà còn cười ha hả: ``Em đang kiểm tra phiên bản phòng khách của nhà mình mà! Phòng khách nhà mình hẹp hơn cả phòng thu này, phải thử xem có đủ chỗ thực hiện động tác không, chứ lên sân khấu lớn có gì mà lo.''

Từ cửa phòng, Naomi vừa đi theo Suzuno vào đón mọi người, thấy Kson tập động tác thì mắt sáng rực. Em bé giơ tay lên cao, giọng trong trẻo vang lên: ``Kson-oneechan làm lại đi ạ! Naomi muốn xem nữa!''

Kson vỗ tay một cái, hào hứng đồng ý ngay: ``Được thôi! Lần này chị làm chuẩn hơn.''

Chị ấy quay lại vị trí ban đầu, hít một hơi sâu như đang chuẩn bị cho một buổi biểu diễn thực thụ. Nhưng khi đến bước xoay người, thay vì bước sang bên phải như tôi đã làm, cô lại bước ngược sang bên trái, cả người xoay tròn nhẹ nhàng nhưng hoàn toàn sai hướng.

Tôi bật cười thành tiếng, tay chỉ vào cô: ``Chị vừa xoay nhầm bên rồi đó! Phải là bên phải chứ không phải trái đâu!''

Kson dừng lại, giả vờ nghiêm túc chỉnh lại cổ áo, như thể mọi thứ đều nằm trong kế hoạch. ``Đó là một biến thể mới! Chị sáng tạo ra đấy, độc nhất vô nhị luôn!''

Ngay lúc đó, loa trên trần nhà bật lên, giọng đều đặn của Game vang ra: ``Tôi đã phân tích dữ liệu chuyển động của Kson. Xác nhận đây là một động tác hoàn toàn khác, không trùng khớp với bất kỳ khung hình nào của chuỗi động tác gốc. Tôi có thể lưu lại biến thể này vào cơ sở dữ liệu nếu cần.''

Cả phòng cùng bật cười, Kson quay người chỉ vào chiếc đồng hồ treo tường, giọng giả bộ tức giận: ``Ai hỏi ông đâu! Ông cứ lo công việc của ông đi, đừng có bình luận động tác của chị.''

Sau khi cười đùa xong, mọi người cùng thu dọn đồ đạc còn lại, bước ra khỏi phòng thu và tắt đèn phía sau. Khoảng trống trên sàn nhà nơi tôi đã tập luyện hàng tuần giờ im lặng, nhưng những tiếng cười nói vang vọng vẫn còn như giữ lại những khoảnh khắc đáng nhớ của buổi tập cuối cùng. Cuối chuyến hành trình chuẩn bị, tất cả mọi người đều sẵn sàng cho ngày trình diễn chính thức, mang theo trong lòng những bài học quý giá và niềm vui của những ngày cùng nhau cố gắng.

\end{document}